\documentclass[11pt]{article}
\usepackage{preamble}
\setlength{\parskip}{4pt}

\begin{document}

\daybanner{AP Statistics \textperiodcentered\ Topic 3.5 (CED 1.13) \textperiodcentered\ B: 2026-10-02 \textperiodcentered\ E: 2026-10-14 \textperiodcentered\ TEACHER KEY}{Which Design Supports Cause?}

\sectionbanner{CED TETHER}

{\small
\begin{itemize}[leftmargin=1.5em,itemsep=2pt,topsep=2pt]
  \item \textbf{Skill 1.C} --- Describe an appropriate method for gathering and representing data.
  \item \textbf{Skill 1.B} --- Identify key and relevant information to answer a question or solve a problem.
  \item \textbf{LO VAR-3.A} --- Identify the components of an experiment. [Skill 1.C]
  \item \textbf{EK VAR-3.A.1} --- The experimental units are the individuals (which may be people or other objects of study) that are assigned treatments.
  \item \textbf{EK VAR-3.A.2} --- An explanatory variable (or factor) in an experiment is a variable whose levels are manipulated intentionally. The levels or combination of levels are called treatments.
  \item \textbf{EK VAR-3.A.3} --- A response variable in an experiment is an outcome from the experimental units that is measured after the treatments have been administered.
  \item \textbf{LO VAR-3.B} --- Describe elements of a well-designed experiment. [Skill 1.B]
  \item \textbf{EK VAR-3.B.1} --- A well-designed experiment includes comparison, random assignment, replication, and control of potential confounding variables where appropriate.
  \item \textbf{EK VAR-3.C.1} --- Random assignment tends to balance uncontrolled variables so that differences in responses can be attributed to treatments.
  \item \textbf{EK VAR-3.C.4} --- In a double-blind experiment neither the subjects nor the research team members who interact with them know which treatment a subject is receiving.
  \item \textbf{EK VAR-3.C.5} --- A control group receives no treatment of interest or an inactive treatment so the treatment of interest can be evaluated.
\end{itemize}
}
\textbf{Essential question:} Which design features let a researcher attribute a response difference to a treatment?\par
\textbf{DOK-3 skill:} 4.A \textperiodcentered\ \textbf{FRQ pattern:} {\small\texttt{experiment-\allowbreak{}principles-\allowbreak{}which-\allowbreak{}design-\allowbreak{}supports-\allowbreak{}cause}}\par
\textbf{DOK rationale:} Part (c) compares two designs against a causal claim, links random assignment, control, and blinding to specific threats, explains a concrete alternative cause, and bounds generalization from volunteers.\par

\frameworkphaseheader{Do Now (first take) $\rightarrow$ Explore (video + follow-along) $\rightarrow$ Exit (finish + turn in)}{1 $\rightarrow$ 3}{45}{%
\begin{itemize}[leftmargin=1.5em,itemsep=2pt,topsep=2pt]
  \item Board slide up at the bell. Ask for a design choice and one feature.
  \item Begin the video follow-along at minute 5.
  \item Return to the board slide at minute 33. Circulate and separate assignment, control, blinding, and selection.
  \item Collect at the door. Score part (c) only, E/P/I.
\end{itemize}
}{%
\begin{itemize}[leftmargin=1.5em,itemsep=2pt,topsep=2pt]
  \item Choose a design and cite one feature before the lesson.
  \item Complete the video follow-along about experimental units, treatments, responses, and design principles.
  \item Finish (a)--(c), revise the first take in the exit line, and turn in the sheet.
\end{itemize}
}{%
\begin{itemize}[leftmargin=1.5em,itemsep=2pt,topsep=2pt]
  \item What exactly did the researcher manipulate, and what was measured afterward?
  \item Which feature addresses preexisting differences between early and late volunteers?
  \item What does the control group reveal that one treatment group cannot?
  \item Who was randomized, and from what larger population were they not randomly selected?
\end{itemize}
}{Solo teacher. Circulate ~5 pairs. Require each named design feature to be linked to the threat it addresses; do not accept a list of principles without mechanisms.}

\begin{callout}{\IconWarn\ WATCH FOR}{calloutred}
\begin{itemize}[leftmargin=1.5em,itemsep=2pt,topsep=2pt]
  \item Treating signup order or equal group size as random assignment.
  \item Calling the study representative of teenagers because its treatments were randomized.
  \item Naming blinding without saying whose expectations or measurements it protects.
\end{itemize}
\end{callout}

\sectionbanner{SCORING PART (c) --- E / P / I}

\textbf{Expected elements}
\begin{itemize}[leftmargin=1.5em,itemsep=2pt,topsep=2pt]
  \item \textbf{chooses-design-2} (required) --- Chooses Design 2 and identifies coin-flip random assignment as the central causal feature
  \item \textbf{control-and-blinding} (required) --- Explains how the matched control and blinding strengthen the comparison beyond assignment alone
  \item \textbf{signup-order-alternative} (required) --- Names a plausible early-versus-late volunteer difference and links it to reaction time
  \item \textbf{limits-generalization} (required) --- States that volunteers do not automatically represent all teenagers despite random treatment assignment
\end{itemize}

{\small\renewcommand{\arraystretch}{1.25}
\noindent\begin{tabular}{|p{0.5in}|p{5.9in}|}
\hline
\textbf{E} & Chooses Design 2 with random assignment, control, and blinding explained, gives a linked signup-order alternative for Design 1, and limits generalization from volunteers. \\ \hline
\textbf{P} & Chooses Design 2 and names key features but incompletely explains the alternative cause or population limit, or treats blinding as the causal feature while barely addressing assignment. \\ \hline
\textbf{I} & Chooses Design 1 because it has equal group sizes, claims a control alone establishes cause, or generalizes to all teenagers because 30 people participated. \\ \hline
\end{tabular}}

\textbf{Common mistakes}
\begin{itemize}[leftmargin=1.5em,itemsep=2pt,topsep=2pt]
  \item Calling Design 1 randomized because students arrived in an unpredictable order
  \item Treating equal group sizes as a substitute for random assignment
  \item Saying identical cups alone blind participants without the matched drink and concealed assignment
  \item Confusing random assignment among volunteers with random selection of teenagers
\end{itemize}

\clearpage
\focusbanner{TODAY'S PROBLEM --- annotated}

Suppose a researcher recruits 30 student volunteers to study whether an energy drink changes reaction time.\par\smallskip \textbf{Design 1:} The first 15 students to sign up receive the energy drink; the last 15 receive water. A tester who knows each drink measures reaction time.\par \textbf{Design 2:} A coin flip assigns each volunteer to the energy drink or water. The water is flavored to match, both drinks are served in identical opaque cups, and neither the volunteer nor the tester knows which drink was assigned.\par
\begin{center}
\textbf{Two energy-drink designs}\par\smallskip
\begin{tabular}{|l|c|c|c|}
\hline
\textbf{\#} & \textbf{Assign} & \textbf{Control} & \textbf{Blind} \\ \hline
\textbf{1} & Signup & Water & No \\ \hline
\textbf{2} & Coin & Matched & Both \\ \hline
\end{tabular}
\end{center}
\begin{firsttakebox}{FIRST TAKE (5 min)}
Which design would make you more willing to attribute a reaction-time difference to the drink? Name one feature behind your choice.\par
\answer{Not graded. Expect students to notice the coin flip before they can explain why it matters.}
\end{firsttakebox}

\textit{Rules callout ``PRINCIPLES OF EXPERIMENTAL DESIGN (keep on the page)'' is printed on the student sheet.}\par\medskip

\dokbadge{1}~\textbf{(a)} Identify the experimental units, explanatory variable and its levels, and response variable.\par
\answer{Experimental units: the 30 student volunteers. Explanatory variable: assigned drink, with levels energy drink and matched water. Response variable: measured reaction time, with its recorded time unit.}

\dokbadge{2}~\textbf{(b)} For each design, state whether it has random assignment, a control group, and blinding. Explain one job performed by each feature.\par
\answer{Design 1 has a water control and replication but no random assignment or blinding. Design 2 has coin-flip random assignment, a matched water control, replication, and double blinding. Random assignment tends to balance lurking variables; the control supplies a comparison; blinding reduces expectancy and measurement effects.}

\dokbadge{3}~\textbf{(c)} Which design could support the claim that the energy drink caused a reaction-time difference? Cite the specific features that settle it, explain what a difference under Design 1 could be due to instead, and decide whether either result could automatically generalize to all teenagers.\par
\answer{Design 2 can support causation for these volunteers because the coin flip randomly assigns treatments, both groups are compared, and blinding limits placebo or tester effects. Under Design 1, early volunteers could differ from late volunteers in motivation, sleep, schedule, or baseline reaction time, so signup order could explain the difference. Neither design automatically generalizes to all teenagers because the 30 participants volunteered and were not randomly selected or shown to represent all teenagers.}

\begin{summaryexitbox}{EXIT}
One thing the video changed about my first take: \hrulefill
\end{summaryexitbox}

\end{document}
